\chapter{Getaltheorie: delers en priemgetallen}\label{ch:g9:arith}

De getaltheorie studeert hele getallen en de manier waarop ze elkaar delen. Haar
hoofdrolspelers zijn de priemgetallen, de bouwstenen waaruit elk geheel getal met
vermenigvuldigingen wordt samengesteld. Het hoofdstuk eindigt met de \omterm{def:g9:arith:gcd}{grootste
gemene deler}, het juiste gereedschap om \omterm{def:g9:fractions:fraction}{breuken} eens en voor altijd te
vereenvoudigen. Dit verhaal gaat, veel verder, door in het bovenbouwvolume en
daarna.

\section{Delers en veelvouden}

\begin{definition}[Deler, veelvoud]\label{def:g9:arith:divisor}
Zij $a$ en $b$ positieve gehele getallen. We zeggen dat $b$ het getal $a$
\emph{deelt}\index{deler} (of dat $b$ een deler van $a$ is, of dat $a$ een
\emph{veelvoud}\index{veelvoud} van $b$ is) wanneer $a = b \times k$ voor een
zeker geheel getal $k$ --- dat wil zeggen wanneer de deling van $a$ door $b$ \omterm{def:g3:division:remainder}{rest}
$0$ laat.
\end{definition}

\begin{example}
De \omterm{def:g9:arith:divisor}{delers} van $24$ zijn $1, 2, 3, 4, 6, 8, 12, 24$ --- ze komen in paren met
\omterm{def:g2:mult:def}{product} $24$: $(1,24)$, $(2,12)$, $(3,8)$, $(4,6)$. De \omterm{def:g5:division:multiple}{veelvouden} van $7$ zijn
$7, 14, 21, 28, \dots$
\end{example}

\begin{proposition}[Deelbaarheidsregels]\label{prop:g9:arith:rules}
Een geheel getal is \omterm{def:g6:wholes:divisible}{deelbaar}:
\begin{itemize}
  \item door $2$ wanneer zijn laatste cijfer \omterm{def:g3:numbers:evenodd}{even} is ($0, 2, 4, 6, 8$);
  \item door $5$ wanneer zijn laatste cijfer $0$ of $5$ is;
  \item door $10$ wanneer zijn laatste cijfer $0$ is;
  \item door $3$ (respectievelijk $9$) wanneer de \omterm{def:g1:addition:def}{som} van zijn cijfers \omterm{def:g6:wholes:divisible}{deelbaar} is
        door $3$ (respectievelijk $9$);
  \item door $4$ wanneer zijn laatste twee cijfers een getal vormen dat \omterm{def:g6:wholes:divisible}{deelbaar}
        is door $4$.
\end{itemize}
\end{proposition}
\admitted

\begin{example}
$7\,215$ eindigt op $5$: \omterm{def:g6:wholes:divisible}{deelbaar} door $5$. Zijn cijfersom is
$7 + 2 + 1 + 5 = 15$, \omterm{def:g6:wholes:divisible}{deelbaar} door $3$ maar niet door $9$: dus is $7\,215$
\omterm{def:g6:wholes:divisible}{deelbaar} door $3$, niet door $9$. En inderdaad is
$7\,215 = 3 \times 5 \times 481$.
\end{example}

\section{Priemgetallen}

\begin{definition}[Priemgetal]\label{def:g9:arith:prime}
Een \emph{priemgetal}\index{priemgetal} is een geheel getal $\geq 2$ waarvan de
enige \omterm{def:g9:arith:divisor}{delers} $1$ en het getal zelf zijn. De priemgetallen onder $30$ zijn
\[
2,\ 3,\ 5,\ 7,\ 11,\ 13,\ 17,\ 19,\ 23,\ 29 .
\]
Het getal $1$ is \emph{geen} priemgetal (per afspraak), en een geheel getal
$\geq 2$ dat niet priem is, noem je \emph{samengesteld}.
\end{definition}

\begin{theorem}[Priemfactorontbinding]\label{thm:g9:arith:factorization}
Elk geheel getal $\geq 2$ is een \omterm{def:g2:mult:def}{product} van priemgetallen, en die ontbinding is
op de orde van de factoren na uniek.
\end{theorem}
\admitted

\begin{method}[Een geheel getal ontbinden]\label{met:g9:arith:factorization}
Deel herhaaldelijk door het kleinst mogelijke \omterm{def:g9:arith:prime}{priemgetal}, tot je bij $1$ uitkomt:
\begin{enumerate}
  \item probeer $2$ zolang het getal \omterm{def:g3:numbers:evenodd}{even} is;
  \item probeer daarna $3$, dan $5$, dan $7$, \dots\ (alleen priemgetallen);
  \item stop wanneer het \omterm{def:g3:division:remainder}{quotiënt} $1$ is; verzamel de factoren met hun
        \omterm{def:g8:powers:def}{exponenten}.
\end{enumerate}
Het is genoeg de priemgetallen $p$ te proberen waarvoor $p^2$ het huidige getal
niet overtreft: \omterm{def:g9:arith:divisor}{deelt} geen van hen het getal, dan is het getal zelf priem.
\end{method}

\begin{example}\label{ex:g9:arith:factorization}
Ontbind $360$, één deling per keer:
\[
360 = 2 \times 180, \quad
180 = 2 \times 90, \quad
90 = 2 \times 45, \quad
45 = 3 \times 15, \quad
15 = 3 \times 5,
\]
en dus
\[
360 = 2 \times 2 \times 2 \times 3 \times 3 \times 5 = 2^3 \times 3^2
\times 5 .
\]
\end{example}

\begin{omfigure}
\begin{tikzpicture}[scale=0.78]
  \node (n360) at (0,0) {$360$};
  \node (n180) at (1.2,-1) {$180$};
  \node (n90) at (2.4,-2) {$90$};
  \node (n45) at (3.6,-3) {$45$};
  \node (n15) at (4.8,-4) {$15$};
  \node (n5) at (6.0,-5) {$5$};
  \node[omThm] (p2a) at (-1.2,-1) {$2$};
  \node[omThm] (p2b) at (0,-2) {$2$};
  \node[omThm] (p2c) at (1.2,-3) {$2$};
  \node[omThm] (p3a) at (2.4,-4) {$3$};
  \node[omThm] (p3b) at (3.6,-5) {$3$};
  \draw (n360) -- (n180); \draw (n360) -- (p2a);
  \draw (n180) -- (n90); \draw (n180) -- (p2b);
  \draw (n90) -- (n45); \draw (n90) -- (p2c);
  \draw (n45) -- (n15); \draw (n45) -- (p3a);
  \draw (n15) -- (n5); \draw (n15) -- (p3b);
\end{tikzpicture}

{\small De factorboom van $360$: elke stap splitst de kleinste priemfactor af (in
rood). Lees je de rode bladeren en de laatste $5$, dan staat er
$360 = 2^3 \times 3^2 \times 5$.}
\end{omfigure}

\begin{theorem}[Euclides]\label{thm:g9:arith:euclid}
Er zijn oneindig veel priemgetallen.
\end{theorem}

\begin{proof}
Stel dat er maar een eindig aantal zouden zijn, zeg $p_1, p_2, \dots, p_k$, en
bekijk
\[
N = p_1 \times p_2 \times \dots \times p_k + 1 .
\]
$N$ door een willekeurige $p_i$ delen laat \omterm{def:g3:division:remainder}{rest} $1$, dus \omterm{def:g9:arith:divisor}{deelt} geen enkele $p_i$
het getal $N$. Maar $N \geq 2$ heeft minstens één priemdeler
(\cref{thm:g9:arith:factorization}) --- een \omterm{def:g9:arith:prime}{priemgetal} dat niet in onze lijst
staat. Tegenspraak: geen enkele eindige lijst kan alle priemgetallen bevatten.
\end{proof}

\section{De grootste gemene deler}

\begin{definition}[ggd]\label{def:g9:arith:gcd}
De \emph{grootste gemene deler}\index{grootste gemene deler} van twee positieve
gehele getallen $a$ en $b$, geschreven $\gcd(a, b)$, is het grootste getal dat
beide \omterm{def:g9:arith:divisor}{deelt}. Wanneer $\gcd(a, b) = 1$, heten de getallen \emph{relatief
priem}\index{relatief priem}: ze hebben geen enkele \omterm{def:g9:arith:divisor}{deler} gemeen behalve $1$.
\end{definition}

\begin{example}
\omterm{def:g9:arith:divisor}{Delers} van $18$: $1, 2, 3, 6, 9, 18$. \omterm{def:g9:arith:divisor}{Delers} van $24$: $1, 2, 3, 4, 6, 8, 12,
24$. Gemeenschappelijke \omterm{def:g9:arith:divisor}{delers}: $1, 2, 3, 6$; dus is $\gcd(18, 24) = 6$. De
getallen $15$ en $28$ zijn \omterm{def:g9:arith:gcd}{relatief priem}.
\end{example}

\begin{proposition}[ggd uit de ontbindingen]\label{prop:g9:arith:gcdfact}
De ggd van twee gehele getallen is het \omterm{def:g2:mult:def}{product} van de priemgetallen die in
\emph{beide} ontbindingen voorkomen, elk genomen met de \emph{kleinste} van zijn
twee \omterm{def:g8:powers:def}{exponenten}.
\end{proposition}
\admitted

\begin{example}\label{ex:g9:arith:gcdfact}
$360 = 2^3 \times 3^2 \times 5$ en $84 = 2^2 \times 3 \times 7$.
Gemeenschappelijke priemgetallen: $2$ (\omterm{def:g8:powers:def}{exponenten} $3$ en $2$: houd $2$) en $3$
(\omterm{def:g8:powers:def}{exponenten} $2$ en $1$: houd $1$). Dus is
\[
\gcd(360, 84) = 2^2 \times 3 = 12 .
\]
\end{example}

\begin{theorem}[Algoritme van Euclides]\label{thm:g9:arith:euclidalgo}
Is $a = bq + r$ de deling van $a$ door $b$ met \omterm{def:g3:division:remainder}{rest} $r$, dan geldt
\[
\gcd(a, b) = \gcd(b, r).
\]
Herhaal je de delingen tot de \omterm{def:g3:division:remainder}{rest} $0$ is, dan is de ggd van $a$ en $b$ de
\emph{laatste \omterm{def:g3:division:remainder}{rest} verschillend van \omterm{def:g1:counting:zero}{nul}}.
\end{theorem}

\begin{proof}
Uit $a = bq + r$: elk getal dat $b$ en $r$ \omterm{def:g9:arith:divisor}{deelt}, \omterm{def:g9:arith:divisor}{deelt} ook $bq + r = a$; en uit
$r = a - bq$: elk getal dat $a$ en $b$ \omterm{def:g9:arith:divisor}{deelt}, \omterm{def:g9:arith:divisor}{deelt} ook $r$. De paren $(a, b)$ en
$(b, r)$ hebben dus precies dezelfde gemeenschappelijke \omterm{def:g9:arith:divisor}{delers} --- en in het
bijzonder dezelfde grootste. Omdat de \omterm{def:g3:division:remainder}{resten} strikt dalen, stopt het algoritme, en
$\gcd(x, 0) = x$ levert de laatste \omterm{def:g3:division:remainder}{rest} verschillend van \omterm{def:g1:counting:zero}{nul}.
\end{proof}

\begin{example}\label{ex:g9:arith:euclidalgo}
Bereken $\gcd(1071, 462)$:
\begin{align*}
1071 &= 462 \times 2 + 147, \\
462 &= 147 \times 3 + 21, \\
147 &= 21 \times 7 + 0 .
\end{align*}
De laatste \omterm{def:g3:division:remainder}{rest} verschillend van \omterm{def:g1:counting:zero}{nul} is $21$: $\gcd(1071, 462) = 21$.
\end{example}

\begin{method}[Een breuk volledig vereenvoudigen]\label{met:g9:arith:simplify}
Om $\dfrac ab$ in eenvoudigste vorm te schrijven:
\begin{enumerate}
  \item bereken $d = \gcd(a, b)$, bijvoorbeeld met het algoritme van Euclides;
  \item deel \omterm{def:g4:fractions:def}{teller} en \omterm{def:g4:fractions:def}{noemer} door $d$: $\dfrac ab = \dfrac{a \div d}{b \div d}$;
  \item de \omterm{def:g9:fractions:fraction}{breuk} die je zo krijgt is \emph{onvereenvoudigbaar}: haar \omterm{def:g4:fractions:def}{teller} en
        \omterm{def:g4:fractions:def}{noemer} zijn \omterm{def:g9:arith:gcd}{relatief priem}.
\end{enumerate}
\end{method}

\begin{example}
$\dfrac{462}{1071} = \dfrac{462 \div 21}{1071 \div 21} = \dfrac{22}{51}$, en
$\gcd(22, 51) = 1$: onvereenvoudigbaar.
\end{example}

\section{Oefeningen}

\begin{exercise}[$\star$]\label{exo:g9:arith:1}
Noem alle \omterm{def:g9:arith:divisor}{delers} van $36$, van $45$ en van $17$.
\end{exercise}

\begin{exercise}[$\star$]\label{exo:g9:arith:2}
Bepaal met de deelbaarheidsregels of $2\,346$ \omterm{def:g6:wholes:divisible}{deelbaar} is door $2$, door $3$, door
$4$, door $5$ en door $9$.
\end{exercise}

\begin{exercise}[$\star$]\label{exo:g9:arith:3}
Geef de priemfactorontbinding van $72$, $150$, $210$ en $121$.
\end{exercise}

\begin{exercise}[$\star$]\label{exo:g9:arith:4}
Is $101$ priem? En $91$? En $143$? Verantwoord met de stopregel van
\cref{met:g9:arith:factorization}.
\end{exercise}

\begin{exercise}[$\star$]\label{exo:g9:arith:5}
Bereken $\gcd(48, 60)$ op twee manieren: door de gemeenschappelijke \omterm{def:g9:arith:divisor}{delers} op te
noemen, en uit de priemfactorontbindingen.
\end{exercise}

\begin{exercise}[$\star\star$]\label{exo:g9:arith:6}
Bereken met het algoritme van Euclides $\gcd(255, 154)$, en daarna
$\gcd(1053, 325)$. Schrijf elke deling uit.
\end{exercise}

\begin{exercise}[$\star\star$]\label{exo:g9:arith:7}
Maak de \omterm{def:g9:fractions:fraction}{breuk} $\dfrac{588}{504}$ onvereenvoudigbaar. (Bereken de ggd met de
methode van je keuze, en deel dan.)
\end{exercise}

\begin{exercise}[$\star\star$]\label{exo:g9:arith:8}
Een bloemist heeft $84$ rozen en $126$ tulpen en wil er identieke boeketten van
maken, met alle bloemen op en zo veel boeketten mogelijk. Hoeveel boeketten kan
ze maken, en wat zit er in elk?
\end{exercise}

\begin{exercise}[$\star\star$]\label{exo:g9:arith:9}
Twee veerboten vertrekken om 8:00 van dezelfde kade. De ene vaart elke $24$
minuten uit, de andere elke $36$ minuten. Hoe laat vertrekken ze de volgende keer
samen? (Zoek het kleinste gemeenschappelijke \omterm{def:g9:arith:divisor}{veelvoud} van $24$ en $36$; de
ontbindingen helpen.)
\end{exercise}

\begin{exercise}[$\star\star\star$]\label{exo:g9:arith:10}
Zij $n$ een positief geheel getal.
\begin{enumerate}
  \item Toon aan dat $\gcd(n, n+1) = 1$ (opeenvolgende gehele getallen zijn altijd
        \omterm{def:g9:arith:gcd}{relatief priem}).
  \item Leid af dat de \omterm{def:g9:fractions:fraction}{breuk} $\dfrac{n}{n+1}$ altijd onvereenvoudigbaar is.
\end{enumerate}
\end{exercise}

\section{Opgave: waterkannen, cicaden en honderd kluisjes}

\begin{problem}[{Weekendopgave --- de ggd beslist welke hoeveelheden twee kannen
kunnen afmeten; priemgetallen beschermen cicaden; en de kluisjes die open blijven
zijn de volkomen kwadraten}]
\label{pb:g9:arith:1}
Drie puzzels die op raadsels lijken en in werkelijkheid getaltheorie zijn: water
afmeten met kannen zonder maatstreep (de ggd in vermomming), levenscycli van
insecten die tot priemgetallen zijn geëvolueerd, en een beroemde gang met honderd
kluisjes waarvan de eindstand door het tellen van \omterm{def:g9:arith:divisor}{delers} wordt beslist. Alles
draait op de machinerie van dit hoofdstuk: deelbaarheid, priemfactorontbinding
(\cref{thm:g9:arith:factorization}) en het algoritme van Euclides
(\cref{thm:g9:arith:euclidalgo}).

\medskip
\noindent\textbf{Deel I --- De waterkannen.}
Je staat bij een fontein met twee kannen zonder maatstreep, van $5$ L en $3$ L.
Toegestane zetten: een kan tot de rand vullen, een kan volledig leegmaken, de ene
kan in de andere gieten tot de bron leeg of het doel vol is.
\begin{enumerate}
  \item Meet precies $1$ L af. (Beschrijf je reeks zetten en de \omterm{def:g2:measure:units}{inhoud} van de twee
        kannen na elke zet.)
  \item Meet precies $4$ L af --- de puzzel uit een beroemde actiefilm. (Het kan
        in zes zetten.)
  \item Welke hele aantallen liter van $1$ tot $8$ kun je voorleggen (in één kan,
        of verdeeld over de twee)? Vul de lijst aan en hergebruik je reeksen.
  \item Nieuwe kannen: $6$ L en $4$ L. Probeer $1$ L af te meten --- en leg daarna
        uit waarom het hopeloos is: ga na dat elk van de drie toegestane zetten de
        \omterm{def:g2:measure:units}{inhoud} van elke kan een \omterm{def:g9:arith:divisor}{veelvoud} van $2$ laat, zodat elke bereikbare
        hoeveelheid \omterm{def:g3:numbers:evenodd}{even} is.
  \item Het argument van vraag 4 werkt in het algemeen: met kannen van $a$ en $b$
        liter is elke bereikbare hoeveelheid een \omterm{def:g9:arith:divisor}{veelvoud} van $\gcd(a, b)$.
        Bereken $\gcd(6, 4)$ en $\gcd(5, 3)$, en zeg wat de wet voor elk paar
        kannen voorspelt.
\end{enumerate}

\medskip
\noindent\textbf{Deel II --- Euclides bij de fontein.}
\begin{enumerate}[resume]
  \item Bereken met het algoritme van Euclides: $\gcd(91, 65)$ en
        $\gcd(2\,026, 46)$.
  \item Leg in je eigen woorden uit waarom de hoeveelheden die in de kannen
        opduiken, de \omterm{def:g3:division:remainder}{resten} van Euclides in vermomming zijn: vul met kannen van
        $13$ L en $5$ L herhaaldelijk de kleine kan en giet die in de grote leeg
        (en maak de grote kan leeg zodra hij vol is). Welke nieuwe hoeveelheden
        duiken het eerst op --- en vergelijk ze met de \omterm{def:g3:division:remainder}{resten} in het algoritme van
        Euclides voor $(13, 5)$.
  \item Leid het antwoord van de kampioen af: kun je met kannen van $13$ en $5$
        liter precies $1$ L afmeten? Verantwoord het in één regel met vraag 5 en
        $\gcd(13, 5)$.
  \item Een snel bewijs van relatieve priemheid in de stijl van
        \cref{exo:g9:arith:10}: toon aan dat $\gcd(n, 2n + 1) = 1$ voor elk
        positief geheel getal $n$. (Wat moet een gemeenschappelijke \omterm{def:g9:arith:divisor}{deler} van $n$
        en $2n + 1$ delen?)
  \item Twee bussen vertrekken om 7:00 samen van de eindhalte; de ene rijdt elke
        $12$ minuten uit, de andere elke $18$. Noem de volgende vertrektijden van
        elk en zoek het eerste ogenblik waarop ze weer samen vertrekken. Ga op dit
        voorbeeld de mooie wet na: (eerste gemeenschappelijke \omterm{def:g9:arith:divisor}{veelvoud}) $\times$
        ggd $=$ \omterm{def:g2:mult:def}{product} van de twee getallen --- en toets ze nog eens op $5$ en
        $3$.
\end{enumerate}

\medskip
\noindent\textbf{Deel III --- Cicaden, \omterm{def:g9:arith:divisor}{delers} en kluisjes.}
\begin{enumerate}[resume]
  \item Bepaalde Noord-Amerikaanse cicaden komen slechts elke $17$ jaar boven;
        stel dat de populatie van een roofdier elke $4$ jaar een piek kent.
        Gebeuren beide dit jaar, over hoeveel jaar valt een opkomst dan voor het
        eerst weer samen met een piek? Dezelfde vraag als de cyclus van de cicaden
        $16$ jaar zou zijn --- hoe vaak zouden ze dan afgeslacht worden? Leg in één
        zin uit waarom de evolutie de cyclus naar een \emph{priem} lengte duwde.
  \item Tel met de ontbinding $360 = 2^3 \times 3^2 \times 5$ de \omterm{def:g9:arith:divisor}{delers} van $360$
        zonder ze op te noemen: een \omterm{def:g9:arith:divisor}{deler} kiest een \omterm{def:g8:powers:def}{exponent} voor $2$ (vier
        keuzen: $0, 1, 2, 3$), een voor $3$ en een voor $5$. Hoeveel \omterm{def:g9:arith:divisor}{delers} zijn er
        in totaal?
  \item Toon aan dat in de ontbinding van een volkomen kwadraat $n = m^2$ elk
        \omterm{def:g9:arith:prime}{priemgetal} een \emph{\omterm{def:g3:numbers:evenodd}{even}} \omterm{def:g8:powers:def}{exponent} draagt. Leid af, zonder één
        vierkantswortel te berekenen, dat $360$ geen volkomen kwadraat is.
  \item Koppel elke \omterm{def:g9:arith:divisor}{deler} $d$ van $n$ aan zijn partner $\frac{n}{d}$ (voor
        $n = 36$: $1 \leftrightarrow 36$, $2 \leftrightarrow 18$,
        $3 \leftrightarrow 12$, $4 \leftrightarrow 9$, $6 \leftrightarrow 6$).
        Wanneer is een \omterm{def:g9:arith:divisor}{deler} zijn eigen partner? Leid het criterium af: $n$ heeft
        een \emph{\omterm{def:g3:numbers:evenodd}{oneven}} aantal \omterm{def:g9:arith:divisor}{delers} precies dan als $n$ een volkomen kwadraat
        is. Ga het na op $36$ en op $360$.
  \item De honderd kluisjes. De kluisjes $1$ tot $100$ zijn in het begin dicht.
        Leerling $1$ verandert de stand van elk kluisje; leerling $2$ verandert de
        stand van de kluisjes $2, 4, 6, \dots$; leerling $k$ verandert de stand van
        de \omterm{def:g5:division:multiple}{veelvouden} van $k$; en zo verder tot leerling $100$. Leg uit welke
        leerlingen kluisje $n$ aanraken, hoeveel keer de stand ervan verandert, en
        --- met vraag 14 --- precies welke kluisjes op het einde open staan.
        Hoeveel staan er open?
\end{enumerate}
\end{problem}
